关于"哪种体位更容易怀孕"，民间流传着不少说法：男上位、后入式、射精后垫高臀部、甚至"倒立"。这些说法的\textbf{科学证据并不充分}。

\noindent\textbf{体位与受孕的真实关系}：
\begin{itemize}
	\item 目前\textbf{没有可靠证据}表明某一种体位能显著提高妊娠率。射精后精子会迅速进入宫颈黏液，其上行主要依赖精子自身的运动能力、宫颈黏液的质量与排卵时机的配合，而非重力；
	\item 男上位、后入式等插入较深的体位，理论上可使精液更靠近宫颈口，射精后平躺或屈膝仰卧 10--15 分钟，可能\textbf{略微}减少精液外流——但尚无高质量研究证明这能提升妊娠率；
	\item "射精后倒立"既无必要，也有受伤风险，不建议尝试。
\end{itemize}

\noindent\textbf{真正有助于怀孕的要点}：
\begin{itemize}
	\item \textbf{抓住排卵期}：排卵前 1--2 天至排卵日同房受孕概率最高，可结合基础体温、排卵试纸或周期推算把握；
	\item \textbf{规律而非"攒着"}：备孕期间每 1--2 天同房一次即可，长期禁欲反而可能降低精子质量；
	\item \textbf{慎选润滑剂}：部分润滑剂会影响精子活力，备孕时优先选择明确标"精子友好"的产品，或改用足量前戏；
	\item \textbf{健康为本}：戒烟限酒、控制体重、规律作息、均衡营养，对男女双方的生育力都有帮助；女方孕前补充叶酸。
\end{itemize}

\begin{tcolorbox}[colback=pink!5!white,colframe=red!50!black,title=贴心小叮咛]
若规律无避孕同房已满 1 年仍未怀孕（女方 35 岁以上为 6 个月），应夫妻双方一同到生殖医学科就诊评估，而不是把希望寄托在"换个姿势"上。
\end{tcolorbox}
